\section*{Additional Course Policies}
\subsection*{Late Assignment/Incomplete Assignment Policy:}
In general, you should plan ahead and submit assignments on time or early if you are able. 
Falling behind can cause many problems in this course. 
However, we understand that scheduling can be difficult at times. 
\textbf{You may request a 24-hour extension} 
(longer extensions are not generally considered due to the pace of the course) 
on an assignment using the \textbf{Extension Request Form} 
found just under the colorful Airtable Calendar at the top of the course Moodle site. 
If you do request an extension, 
you should do your best to request the extension 
before noon (12 p.m.) Eastern time of the day the assignment is due.  
There is no hard limit to how many extensions you may request, 
but we reserve the right to deny what appear to be unreasonable requests, 
or requests in which we believe granting the extension 
would cause more harm than good. 
Please remember that we expect students to act professionally. 
It can help to ask yourself whether you would 
make such a request to a manager while at work 
and whether the request is based on reasonable circumstances. 

Here are some general guidelines:
\begin{itemize}
\itemsep0em
\item Send extension requests as soon as feasible, 
generally before noon (12:00 p.m.) on the date the assignment is due. 
If you are very ill, the timing is more flexible, but the sooner the better. 
\item It is OK to request an extension and then not use it. 
    % If you do not use a requested extension, 
    % it does not ``count'' as a used extension. 
\item Putting the assignment off until the last moment 
    (such as starting 4 hours before the deadline) 
    is NOT a justification for an extension.  
    We reserve the right to see your work-in-progress. 
\item \textbf{NOTE:} Most assignments will usually take more work 
    than just sitting down and coding everything in an hour. 
    You should plan accordingly and start your assignments early. 
    The course Airtable calendar indicates when you should start looking at each assignment. 
\item The typical extension is 24 hours. 
    On very rare occasions, we may allow 48 hours. 
\item We keep track of extensions per student.  
    If you seem to be overusing the privilege,  
    we may stop allowing extensions or make special arrangements for you. 
\item If the reason your work is late is because 
    you are having trouble completing the assignment 
    (e.g., a bug you can't fix or just difficulty understanding), 
    reach out for help (see ``How to Get Help''). 
    Don't just assume that more time will fix the problem. 
\item Please understand that the schedule of assignments is very tight. 
    Once you have completed one assignment, 
    it is usually then time to start the next as soon as possible. 
    So, if you request an extension on an assignment, 
    please realize that extra time will now take time away from your next assignment. 
    Don't let this snowball, 
    causing every assignment to be later than the last. 
\end{itemize}

\subsection*{Attendance Policy:}
Attendance in class is an essential part of your learning. 
We trust that you will attend every class session unless 
an illness, athletic/extracurricular event, etc., prevents you from doing so. 
When you must miss class, please let me know as early as possible. 
After \textbf{four unexcused absences}, 
your final grade may be reduced by \textbf{one full letter grade} for each additional unexcused absence.
See the \href{https://catalog.rose-hulman.edu/rules-procedures/attendance/}{Registrar's page} 
for additional information on institute-wide attendance policies. 


\subsection*{Course Participation Policy:}
When you are physically present in class, 
We expect you to also be \emph{mentally} present and focused on CSSE 220 activities. 
If you and/or your group finishes an activity early, 
you can work on assignments, review/organize your notes, or 
take other actions that further your CSSE 220 learning. 
We know it may be tempting to hop over to an assignment for a different course, 
but context-switching hinders your learning 
(and \href{https://dl.acm.org/doi/pdf/10.1145/3084100.3084116}{harms productivity of software developers}). 


\subsection*{Course Academic Integrity Policy:}
We believe that everyone in this course is fundamentally honest, 
and we will help you learn the conventions of academic integrity,  
such as citing sources when you use others' ideas. 
It is essential for all students to cite any and all sources of help 
received in completing coursework. 
This practice not only fosters a culture of honesty and transparency but also 
prevents misunderstandings that might otherwise escalate to formal proceedings. 
Students should also be aware of what is appropriate help on homework assignments---%
see 
\href{https://rosehulman.sharepoint.com/:w:/r/sites/CS/Shared%20Documents/Pilot%20Integrity%20Policy/IntegrityPilotPolicy%20-%20What%20Constitutes%20Misconduct.docx?d=w72d3dc4b233d4df68e10736359a06c80&csf=1&web=1&e=QWvFcM}
{What Constitutes Misconduct}. 
To ensure fairness and responsibility, 
any instances of suspected misconduct will be handled through the CSSE Integrity Committee. 

If a case of suspected misconduct arises, 
it will be submitted to the CSSE Integrity Committee for review 
(see 
\href{https://rosehulman.sharepoint.com/:w:/r/sites/CS/Shared%20Documents/Integrity%20Policy/Integrity%20Policy%20-%20Procedures%20-%20rev%2004-27-2026.docx?d=w919879b57499452894ee53ed1a267eac&csf=1&web=1&e=rnHf8y}{Integrity Policy---Procedures} 
and possible penalties 
\href{https://rosehulman.sharepoint.com/:w:/s/CS/IQA8KW6RGqOTRL291reO0tPdAfY1xviT03qG5Tcyx1hu6Ok?e=ewVAIx}{IntegrityPilotPolicy---Penalties \& Evidence}) 
Exceptions may be made for cases which the instructor deems minor 
or which arise late in the term. 
The instructor may choose to handle such cases directly 
(without CSSE Integrity Committee involvement). 
The process includes an initial review of the evidence by the committee, 
a time for students to explain or admit to potential misconduct, and 
potentially a hearing to examine the circumstances and evidence. 
Students are encouraged to continue their studies and 
engage with the course material and instructor normally 
throughout the investigation. 

\subsection*{Generative AI:}
% Suggested policy regarding generative AI, from the Rose-Hulman’s Culture of Integrity Implementation Committee: 
It is expected that any work submitted for assessment represents 
the intellectual work of the individual(s) submitting the work. 
Any attempt to pass off the intellectual work of another 
(including the work generated by Large Language Models (LLMs) like ChatGPT) 
as their own or without proper attribution is an example of academic misconduct. 
\textbf{When assignments \emph{do} allow/require 
the use of LLMs/genAI tools, 
clearly indicate which part(s) it generated, and 
include conversation links.}

